\part{The problem statement and research motivation}

\section{What exactly is the problem?}

\begin{qa}{1}{Define the CARS-SF problem in one precise sentence.}
When several socially critical services share one scarce, failure-prone rural backhaul, decide continuously how
much of that bottleneck each \emph{service} receives (from the operating context and measured demand), enforce
that decision in the data plane, and restore connectivity on its own when the backhaul degrades or fails, while
keeping critical-service SLAs.
\end{qa}

\begin{qa}{2}{What does CARS-SF stand for, and what does each part of the name mean operationally?}
\textbf{C}ontext-\textbf{A}ware \textbf{R}ural \textbf{S}ervice \textbf{S}licing \textbf{F}ramework.
\begin{itemize}
\item \emph{Context-aware}: the priority matrix is chosen by \code{ContextInferenceEngine} from clock, season, alerts
and admin override (\code{context/engine.py}).
\item \emph{Rural}: one shared uplink is the only contention point, with a standby backup backhaul (Fig.~\ref{fig:topo}).
\item \emph{Service slicing}: four classes (healthcare, education, agriculture, general), each with its own HTB class,
floor and SLA (\code{registries/slices.yaml}), not eMBB/URLLC/mMTC.
\item \emph{Framework}: swappable stages behind protocols (\code{core/interfaces.py}), composed from
\code{registries/schedulers.yaml}.
\end{itemize}
The thesis title also says ``autonomous self-healing'': the \code{SelfHealingEngine}.
\end{qa}

\begin{qa}{3}{Is the primary research problem rural connectivity, service-aware allocation, context awareness, slicing, self-healing, or their integration?}
The integration: a single closed loop in which context, measured demand, service slicing and self-healing act on
the \emph{same} bottleneck through shared constraints. The healer constrains the allocator rather than overwriting it
(\code{apply_constraints}). Rural connectivity is the setting, not the contribution.
\honest{Integration is a weaker novelty claim than a new algorithm. Defend it with what integration
\emph{enabled} that the parts alone cannot: backup-capacity ceilings that prevent bufferbloat after failover, and
congestion-aware prediction that prevents false failover (both found and fixed in this work).}
\end{qa}

\begin{qa}{4}{Which is the central contribution, and which are supporting engineering components?}
\textbf{Central:}
\begin{itemize}
\item (C1) service-centric slicing with a context-driven, demand-aware allocator enforced at the rural bottleneck;
\item (C2) a loop-free self-healing state machine with predictive switchover, integrated with allocation;
\item (C3) a packet-matched measurement methodology that makes C1 and C2 evaluable.
\end{itemize}
\textbf{Supporting:} traffic models, harness, matrix runner, 5G deployment.
\textbf{Secondary research results:} SliceNet on measured traffic and the synthetic-to-measured transfer study.
RL is an evaluated alternative with a \emph{negative} result.
\end{qa}

\begin{qa}{5}{What concrete failure or inefficiency exists in a conventional rural network that CARS-SF addresses?}
\begin{itemize}
\item (a) A single FIFO uplink lets bulk downloads and video inflate everyone's queue. In the main study, ``no slicing''
meets the healthcare SLA in \NHcSlaNoSlicing{} of windows vs \NHcSlaProposed{} with slicing.
\item (b) Static per-technology shares (3GPP-style) ignore that a disaster changes \emph{which} service matters.
\item (c) A fiber cut with no automated failover means an outage of the order of the repair time. Without self-healing,
the injected cut cost \NOutCutNoneMean\,s of core-bound outage in each run (until the scripted repair).
\end{itemize}
\end{qa}

\begin{qa}{6}{What do you mean by resilient connectivity? Which measurable properties constitute resilience?}
Six measured properties:
\begin{itemize}
\item time to detect (MTTD) and time to act (control-plane MTTR);
\item service outage: the longest delivery gap of affected flows, from matched packets;
\item packets lost during failover;
\item SLA compliance while degraded;
\item correct localisation (link vs core switch);
\item safe fail-back with no loops or flapping.
\end{itemize}
All six are produced per run (\code{healer.resilience_summary}, \code{analysis/metrics.outage_metrics}).
Availability over months is \emph{not} measured.
\end{qa}

\begin{qa}{7}{What do you mean by context-aware? What qualifies as context, and how does it change a decision?}
Context is information \emph{external} to the traffic that changes social priority:
\begin{itemize}
\item time of day and weekday (school hours);
\item season (harvest months);
\item alerts (flood, fire, medical emergency, fiber cut);
\item an administrative override.
\end{itemize}
It selects the priority matrix (e.g.\ disaster: healthcare weight 10 vs 1), and in timeline runs it also changes the
offered traffic. It does \emph{not} change floors or SLA targets.
\ev{timeline runs switch contexts live; unit tests check the precedence rules (\code{tests/test_context.py}).}{\Lthree}
\end{qa}

\begin{qa}{8}{What do you mean by self-healing? How does it differ from detecting a fault, raising an alert or manual repair?}
Self-healing closes the loop: it detects (link watcher, telemetry, prediction), diagnoses (link vs core switch),
acts on the network (quarantine, backup activation, MAC flush), adapts the allocation (constraints, ceilings), and
reverts safely once the fault has cleared (hold-down + OAM check). An alert stops at detection; manual repair needs a human.
\honest{Post-action verification is partial. Escalation to \textsf{DEGRADED} happens when the core is still unreachable
after failover, but there is no generic ``verify KPI, otherwise roll back'' stage for every action (see Q514).}
\end{qa}

\begin{qa}{9}{What is the exact system boundary of CARS-SF?}
\textbf{Inside:}
\begin{itemize}
\item the control loop (\code{core/}, \code{slicing/}, \code{self_healing/}, \code{scheduler/}, \code{classification/online.py});
\item the qdiscs and filters on the two ends of the rural uplink;
\item the s\_edge-side backhaul ports.
\end{itemize}
\textbf{Outside:}
\begin{itemize}
\item end hosts and their congestion control;
\item the core and edge servers;
\item the radio access network;
\item physical repair.
\end{itemize}
The traffic generators, topology and capture are \emph{experimental apparatus}, not part of the framework.
\end{qa}

\begin{qa}{10}{What does CARS-SF control, what does it observe, and what is outside its control?}
\textbf{Controls:}
\begin{itemize}
\item HTB rates and ceilings per class, in both directions;
\item classification rules (host, app, per-flow);
\item the admin state of the s\_edge backhaul and backup ports;
\item the forwarding database (MAC flush).
\end{itemize}
\textbf{Observes:}
\begin{itemize}
\item interface counters and HTB class stats;
\item RTT probes and link state;
\item OVS stats (on OVS);
\item packet captures (post-run, plus live for the mixed UE).
\end{itemize}
\textbf{Outside its control:} offered load, TCP behaviour, the far end of the backhaul, the radio.
\end{qa}

\begin{qa}{11}{Is CARS-SF a protocol, a management framework, a decision engine, a dataset, an ML model, or a combination?}
Most defensible: \emph{a network-management framework with a closed-loop decision engine}, plus an evaluation
methodology. It is not a protocol and does not change packet formats. Datasets (RSNSD, the measured windows)
and models (SliceNet, PPO) are \emph{artefacts} it produces and consumes.
\end{qa}

\begin{qa}{12}{What is the unit of decision-making?}
The \textbf{service slice} (four HTB classes) for allocation, re-planned every 5\,s. The \textbf{flow} (5-tuple) only
for in-loop classification of a mixed UE. The \textbf{link} for self-healing. There are no per-packet decisions:
the kernel's HTB and fq\_codel handle packets.
\end{qa}

\begin{qa}{13}{What is the unit of failure detection and recovery?}
Detection: a link (carrier transition on \code{s_edge-bh}), reachability per tier (edge vs core RTT probes), and the
per-slice SLA. Recovery: path-level (primary $\leftrightarrow$ backup backhaul) plus slice-level constraints.
\end{qa}

\begin{qa}{14}{What does the system optimise, and what constraints must never be violated?}
It maximises priority-weighted demand satisfaction
$U=\sum_s w_s\min(1,x_s/d_s)/\sum_s w_s$, approximately and greedily (Q432).
\textbf{Hard constraints:}
\begin{itemize}
\item $\sum_s x_s\le C$;
\item $x_s\ge \min(f_s,d_s)$ for the floors $f$ (healthcare 5, education 2, agriculture 1, general 1\,Mb/s);
\item the L2 domain is never looped;
\item learned policies can only rescale priorities, never bypass floors or constraints.
\end{itemize}
\end{qa}

\begin{qa}{15}{Maximise average performance, or guarantee minimum quality for critical applications?}
Minimum guarantees first: floors are allocated before any priority share, and HTB \code{rate} is a reservation.
Above the floors it maximises weighted satisfaction.
\honest{The floors are bandwidth guarantees, not latency guarantees. Latency is protected empirically
(fq\_codel + priority), not provably (Q461).}
\end{qa}

\begin{qa}{16}{Improving average performance vs preserving critical connectivity in an emergency?}
The average can stay flat while the critical service collapses. In the main study, weighted satisfaction is
statistically indistinguishable across schedulers (\NUtilMin--\NUtilMax), yet healthcare SLA compliance differs by
\NSlaDeltaNoSlicing{} between slicing and no slicing. CARS-SF is judged on the second kind of metric.
\end{qa}

\begin{qa}{17}{Which user-facing outcomes improve if CARS-SF works?}
\begin{itemize}
\item Telemedicine and ECG streams keep low delay and loss under load (healthcare p95 $\approx$\NHcPninefiveProposed\,ms in the main study).
\item Classrooms and downloads get the surplus instead of a fixed share.
\item A fiber cut is a sub-second blip instead of a minutes-long outage.
\item A shared multi-service device (a village kiosk) still gets per-service treatment.
\end{itemize}
\end{qa}

\begin{qa}{18}{What is the smallest implementation that could demonstrate the central hypothesis?}
One bottleneck, two classes (critical vs bulk), HTB+fq\_codel, a static context switch and an injected link cut with a
standby path. Measure the critical SLA and outage with and without. The \code{ablation} and \code{resilience} studies
are essentially this, scaled up.
\end{qa}

\begin{qa}{19}{What functionality is unnecessary for proving that hypothesis?}
RL policies, SliceNet, the 5G deployment, the RSNSD generator, PPO calibration and the OVS QoS backend. They answer
\emph{secondary} questions (RQ3, RQ5, RQ6, RQ8). The thesis presents them as such.
\end{qa}

\begin{qa}{20}{What is explicitly outside the scope of the thesis?}
\begin{itemize}
\item radio scheduling and handover;
\item real field deployment;
\item multi-site topologies with several bottlenecks;
\item multipath load balancing;
\item controller high availability;
\item security of context feeds;
\item energy and power-failure modelling;
\item real-trace traffic.
\end{itemize}
\end{qa}

\section{Why rural networks?}

\begin{qa}{21}{Why rural connectivity rather than generic enterprise networking?}
Rural sites typically combine one scarce shared backhaul, little redundancy, long repair times and a few socially
critical services (clinic, school, farm cooperative) whose importance swings with events (harvest, flood). Enterprise
networks usually have redundant links and overprovisioning, so allocation under scarcity and autonomous healing
matter less there.
\honest{This characterisation is argued from the problem framing and literature, not from rural measurements made
in this thesis.}
\end{qa}

\begin{qa}{22}{Which rural characteristics make the problem technically different?}
\begin{itemize}
\item A single contention point carries every service.
\item Demand composition changes with context more than with load.
\item The backhaul is long and fragile, and any backup has lower capacity.
\item Repair takes long, so autonomy matters.
\end{itemize}
\end{qa}

\begin{qa}{23}{Limited bandwidth, intermittent backhaul, power, sparse infrastructure, latency or poor redundancy?}
Represented: limited bandwidth (100\,Mb/s shared, swept 50--150), backhaul failure and degradation, poor redundancy
(a single standby backup at half capacity), and long backhaul delay (20\,ms primary, 40\,ms backup).
Not represented: power instability, sparse radio coverage.
\end{qa}

\begin{qa}{24}{Which constraints are directly represented in the testbed?}
\begin{itemize}
\item bottleneck capacity (HTB root = 100\,Mb/s);
\item propagation delays (netem on destination-facing ports);
\item a backup path of 50\,Mb/s and 40\,ms;
\item failures: fiber cut, gradual degradation (delay +15\,ms and loss +0.3\% per tick), core switch failure.
\end{itemize}
See \code{topology/rural_topology.py} and \code{self_healing/injector.py}.
\end{qa}

\begin{qa}{25}{Which are only assumed or discussed conceptually?}
Power failures, radio coverage gaps, weather-driven radio degradation, the cost of a backup link, multi-village
topologies and actual rural traffic mixes. The radio fields in RSNSD are scenario-injected and excluded from ML.
\end{qa}

\begin{qa}{26}{What evidence supports heterogeneous, competing requirements of rural services?}
The application registry assigns per-application SLAs (e.g.\ ECG 20\,ms, telemedicine 50\,ms, bulk library 300\,ms,
irrigation 500\,ms) and distinct traffic models.
\honest{These are literature-typical values chosen by design (\code{registries/applications.yaml}), not measured
in rural deployments. The heterogeneity is \emph{assumed} and then modelled.}
\end{qa}

\begin{qa}{27}{Rural network vs low-capacity network: interchangeable in your formulation?}
Not quite. A low-capacity network explains the scarcity. ``Rural'' adds a single fragile backhaul with a weaker backup,
socially prioritised services, and context swings. The allocation results would carry over to any single-bottleneck
network; the resilience results depend on the rural assumptions.
\end{qa}

\begin{qa}{28}{Why not solve it by increasing bandwidth?}
\begin{itemize}
\item Capacity is the expensive, slow resource in rural areas.
\item The capacity sweep (\code{sensitivity_capacity}) shows the proposed scheduler's advantage is modest at every capacity.
At 50\,Mb/s in disaster it protects healthcare better (SLA share \NSensCapDisasterProposedFiveZeroHcSla{} vs
\NSensCapDisasterStaticFiveZeroHcSla), without higher utility.
\item More bandwidth does not fix a cut fiber.
\end{itemize}
\end{qa}

\begin{qa}{29}{Why might adding bandwidth fail for intermittent or physical failures?}
A cut link has zero capacity whatever its nominal size. Recovery needs a \emph{different} path, plus fast
detection and safe switching. Bandwidth helps only the backup's quality while on backup.
\end{qa}

\begin{qa}{30}{No alternate route, no spare capacity, no radio coverage?}
Then CARS-SF can only \emph{degrade gracefully}: it detects the fault, gives the surviving capacity (e.g.\ edge-hosted
services) to critical slices, and reports. This is exactly the core-switch-failure case: the outage of core-hosted
services lasts until repair (\NOutCoreFullMean\,s, the same as without healing), while edge services keep running.
\end{qa}

\begin{qa}{31}{Can software restore connectivity when the physical infrastructure is completely unavailable?}
No. Software can only use surviving physical paths faster and more safely, and re-prioritise. The thesis states
this explicitly through the core-switch result (Fig.~\ref{fig:resilience}).
\end{qa}

\begin{qa}{32}{What minimum infrastructure does CARS-SF assume?}
\begin{itemize}
\item one Linux-capable aggregation switch at each end of the bottleneck (tc / HTB);
\item a standby backhaul for healing;
\item the ability to probe edge and core servers;
\item UEs whose addresses map to institutions (or SliceNet for mixed devices).
\end{itemize}
\end{qa}

\begin{qa}{33}{How do you model limited resources without making the network so constrained that your method looks artificially good?}
\begin{itemize}
\item Design demands run from 80\,Mb/s (normal) to 145\,Mb/s (school), i.e.\ 0.8--1.45 of capacity, not extreme overload.
\item Capacity is swept from 50 to 150\,Mb/s and load from 0.5$\times$ to 1.5$\times$.
\item The result is reported honestly: at 0.5$\times$ load the proposed scheduler is \emph{not} better than static
(\NSensLoadSchoolProposedZeropFive{} vs \NSensLoadSchoolStaticZeropFive).
\end{itemize}
\end{qa}

\begin{qa}{34}{How representative are the topologies of actual rural deployments?}
Structurally plausible (access aggregation, a long backhaul, an edge MEC, a core), but parameterised, not surveyed.
It is one topology with one bottleneck. Treat the results as properties of this class of topology.
\end{qa}

\begin{qa}{35}{One backhaul village vs a region with multiple community gateways?}
With multiple gateways you get several bottlenecks and real path choice, so allocation becomes a multi-link problem
(per-path HTB, routing). Healing gains load balancing (ECMP or OpenFlow groups instead of a standby link), and the
controller probably becomes hierarchical (per-village agents). See Q130 and Q156.
\end{qa}

\begin{qa}{36}{Why are rural networks the right setting, not just a motivating example?}
Because the three mechanisms matter most exactly there: scarcity (allocation matters only above 1$\times$ load in
our sweep), context swings (disaster weights), and fragile single backhauls with long repair (self-healing).
\end{qa}

\section{Why these application domains?}

\begin{qa}{37}{Why e-healthcare, e-education and smart agriculture?}
They are the canonical rural public-service pillars, and they span the QoS space: latency-critical low-rate
(ECG, voice), high-throughput elastic (classroom video, library downloads), and sparse periodic IoT (soil, weather).
General internet is the best-effort fourth class.
\end{qa}

\begin{qa}{38}{What are the distinctive QoS requirements of each domain?}
From \code{registries/applications.yaml}:
\begin{itemize}
\item Healthcare: one-way latency 10--50\,ms, loss $\le$1\%, jitter $\le$10\,ms.
\item Education: 150--300\,ms, high throughput.
\item Agriculture: 250--1000\,ms, low rate, loss tolerant.
\item General: 300--500\,ms, best effort.
\end{itemize}
Slice-level targets are in \code{registries/slices.yaml}.
\end{qa}

\begin{qa}{39}{Why is a domain label alone insufficient to determine the priority of every flow?}
A domain contains apps with different needs (ECG 0.08\,Mb/s at 20\,ms vs telemedicine video at 4\,Mb/s with 50\,ms),
and context changes urgency. CARS-SF therefore uses domain $\to$ slice for \emph{bandwidth}, per-app SLAs for
\emph{evaluation}, and fq\_codel \emph{within} a slice so that its sparse flows are not delayed by its bulk flows.
\honest{There is no per-app priority within a slice.}
\end{qa}

\begin{qa}{40}{How does emergency healthcare traffic differ from routine monitoring?}
In the models: emergency voice (talkspurt VoIP) and ambulance on/off bursts vs periodic health-monitoring reports.
Under the \code{medical_emergency} context the healthcare weight rises to 15 and its demand to 95\,Mb/s.
In the network the distinction is by context, not by flow.
\end{qa}

\begin{qa}{41}{Is a live classroom always more important than agricultural telemetry?}
No. The school context weights education 5 vs agriculture 2, the harvest context the opposite (3 vs 5), and floors
keep sparse IoT alive in all cases. Within a slice, fq\_codel keeps sparse telemetry from waiting behind video.
\end{qa}

\begin{qa}{42}{How do you distinguish an application needing throughput from one needing low latency?}
By declared SLA (registry) and by mechanism: throughput comes from HTB rate and surplus, latency from AQM (fq\_codel)
plus HTB priority. The ablation shows that removing borrowing raises school healthcare p95 from
\NAblSchoolFullHc{} to \NAblSchoolFullStrictHc\,ms.
\end{qa}

\begin{qa}{43}{Can two applications in the same domain have conflicting QoS requirements?}
Yes: ECG (strict latency, tiny rate) vs telemedicine video (bursty, larger). They share the healthcare class.
fq\_codel's per-flow queues isolate the sparse ECG flow from video bursts, but bandwidth is shared.
Per-app sub-classes are a possible extension.
\end{qa}

\begin{qa}{44}{How should CARS-SF respond when several critical services compete for the same bottleneck?}
Floors first, then priority-weighted shares capped at demand, then the surplus by priority among unsatisfied slices
(Fig.~\ref{fig:alloc}). If even the floors exceed capacity, each floor is clipped to $C/4$ (Q427).
\end{qa}

\begin{qa}{45}{What policy prevents indefinite starvation of best-effort traffic?}
Three mechanisms:
\begin{itemize}
\item a general-class floor of 1\,Mb/s;
\item work-conserving HTB (best effort borrows idle capacity);
\item healing boosts that expire after 6 ticks.
\end{itemize}
\honest{1\,Mb/s keeps it alive but satisfies little. In disaster runs the general slice's satisfaction is low by design.}
\end{qa}

\begin{qa}{46}{How would you extend the framework to industrial IoT, disaster response or public safety?}
Add a slice entry (class, floor, SLA) in \code{slices.yaml}, applications in \code{applications.yaml}, weights per
context in \code{scenarios.yaml}, and retrain SliceNet. QoS class tables, the allocator and monitoring are data-driven.
\honest{The service list is hard-coded as a 4-tuple in several modules (\code{SERVICE_CLASSES}). A fifth class needs
those constants generalised (Q159).}
\end{qa}

\begin{qa}{47}{Were the domains chosen for real requirements, distinct traffic, or easy classification?}
For requirements and policy relevance. Distinct traffic patterns follow from the applications. Classification was
deliberately made \emph{harder} by excluding ports and DSCP from features. Even so, the traffic models were written
by us, so their separability is partly designed (Q334).
\end{qa}

\section{Research questions and hypotheses}

\begin{qa}{48}{State the primary research question in a testable form.}
On a shared rural bottleneck under the design contexts, does service-centric slicing (context-driven, demand-aware,
with self-healing) achieve higher priority-weighted demand satisfaction and critical-service SLA compliance than
technology-centric (3GPP), static or no slicing? (RQ1, \code{docs/thesis/research_design.md}.)
\end{qa}

\begin{qa}{49}{State the null hypothesis for the main evaluation.}
$H_0$: for each baseline $b$, the paired difference in weighted utility between the proposed scheduler and $b$ has
median zero across (scenario, seed) pairs (Wilcoxon signed-rank).
\end{qa}

\begin{qa}{50}{State the alternative hypothesis.}
$H_1$: the median paired difference is positive (proposed better). It was pre-registered as H1a.
\textbf{Result: not supported.} No baseline difference in weighted utility is significant (smallest $p=\NUtilPMin$).
\end{qa}

\begin{qa}{51}{What is your independent variable?}
The scheduler (seven levels), crossed with context (five levels). In the other studies the IVs are the ablation
variant, capacity, load, UEs per service, healing variant and classification mode.
\end{qa}

\begin{qa}{52}{What are your dependent variables?}
Weighted demand satisfaction, per-slice demand satisfaction, SLA-met share, p95/p99 one-way delay, jitter, loss,
Jain fairness, outage, failover loss, MTTD/MTTR, and macro-F1.
\end{qa}

\begin{qa}{53}{Which variables must be held constant?}
Topology and capacities, traffic models and seeds (identical across schedulers for a given seed), run duration (40\,s),
window length (5\,s), UEs per service (2), host load (4 parallel workers, pacing without busy-waiting) and code version
(the git hash is recorded per run).
\end{qa}

\begin{qa}{54}{Which baseline would falsify the claim that context awareness is useful?}
The \code{no_context} ablation, which uses the normal priority matrix in every context. Context awareness is falsified
if it performs as well as \code{full}. Result: school utility \NAblSchoolNoContextUtil{} vs \NAblSchoolFullUtil{} (full);
disaster \NAblDisasterNoContextUtil{} vs \NAblDisasterFullUtil. \textbf{By this test, context awareness is falsified in our setting}:
demand-capping plus HTB borrowing absorb the effect of the weights.
\honest{Say this plainly. Context still decides \emph{which} slice wins under scarcity, but its measured contribution to outcomes here is nil.}
\end{qa}

\begin{qa}{55}{What result would show self-healing adds nothing beyond static slicing?}
Equal outage with and without healing. Observed: fiber cut \NOutCutFullMean{} vs \NOutCutNoneMean\,s; degradation
\NOutDegFullMean{} vs \NOutDegNoneMean\,s. For a core switch failure, healing indeed adds nothing to the outage
(\NOutCoreFullMean{} vs \NOutCoreNoneMean\,s), and we report that.
\end{qa}

\begin{qa}{56}{What result would show the ML classifier is unnecessary?}
If address-based (host) slicing served the mixed UE as well as in-loop SliceNet. Observed:
\begin{itemize}
\item disaster: equal (\NMixedDisasterHostSla{} vs \NMixedDisasterSlicenetSla);
\item school: SliceNet better (\NMixedSchoolSlicenetSla{} vs \NMixedSchoolHostSla).
\end{itemize}
H5b required a gain in both, so it is \textbf{not supported}. More fundamentally, for single-institution UEs, address
rules suffice and SliceNet \emph{is} unnecessary. Its value is confined to multi-service devices.
\end{qa}

\begin{qa}{57}{What would count as a negative result, and how would you report it?}
Any pre-registered hypothesis whose test fails. The report generator marks verdicts mechanically
(\code{experiments/report.py}). H1a, H1b, H3, H5b, H7 and H9 are reported as \emph{not supported} in
\code{docs/thesis/results.md}, next to the supported ones.
\end{qa}

\begin{qa}{58}{Which claims are hypotheses, which are implementation goals, and which are demonstrated findings?}
\begin{itemize}
\item \textbf{Demonstrated} (L3):
  \begin{itemize}
  \item slicing vs none for healthcare SLA (exploratory, significant);
  \item self-healing outage reduction (H4a, H4b);
  \item SliceNet F1 $\ge$0.9 on held-out runs (H5a);
  \item poor synthetic transfer (H6);
  \item self-healing reduces core-bound outage against no healing (H10, added during the audit, not pre-registered);
  \item 5G-core benefit under contention (H8 is marked supported by its directional criterion, but only the disaster effect is substantive: the school difference is about 1\,ms with $n=3$ and changed sign between the pre- and post-audit runs).
  \end{itemize}
\item \textbf{Not supported}: H1a, H1b, H3, H5b, H7, H9 (context-aware vs context-unaware: +0.002 utility, $p=0.054$).
\item \textbf{Implementation goals only}: OVS QoS backend at scale, load balancing (absent), multi-site.
\end{itemize}
\end{qa}
